\mysection[black]{2. x86lite and LLVM}

\mysubsection{2.1 x86lite}




\BDEF{Calling Convention}{
\textbf{Args 1-6} in \texttt{rdi rsi rdx rcx r8 r9}; args $\ge7$ on stack (pushq in rev. order);
\textbf{return} in \texttt{rax};
\textbf{caller-saves} \texttt{rax rcx rdx rsi rdi r8-r11};
\textbf{callee-saves} \texttt{rbx rbp (top) r12-r15}
}
\textbf{Instr.}: INSTR SRC DEST, reg. with \%, imm. with \$, $(A)$ mem. at the address stored in $A$ \\
\textbf{Sfx.}:b=1,l=2,w=4,q=8 bytes\\
\textbf{Args.}: for $n\ge 7:\text{addr}=((n-7)+2)\cdot 8 + \%\text{rbp}$
\begin{multicols}{4}
\textbf{Prologue}
\begin{lstlisting}
 pushq %rbp
 movq %rsp,%rbp
 subq $x,%rsp
\end{lstlisting}
\columnbreak
\textbf{Epilogue}
\begin{lstlisting}
 addq $x,%rsp
 popq %rbp
 retq
\end{lstlisting}
\columnbreak
\textbf{Push}\\
\begin{lstlisting}
pushq %X:
 rsp <- rsp-8
 Mem[rsp] <- X
\end{lstlisting}
\columnbreak
\textbf{Pop}\\
\begin{lstlisting}
popq %X:
 X <- Mem[rsp]
 rsp <- rsp+8 
\end{lstlisting}
\end{multicols}
\textbf{Flg}: OF = under-/overflow, SF = sign (1=neg.),ZF=0
\textbf{CC}: $e=\text{ZF}$,$ne=\neg\text{ZF}$, $g=\neg \text{ZF}\wedge \text{SF}{=}\text{OF}$,$l=\text{SF}\neq \text{OF}$,$ge=\text{SF}=\text{OF}$,$le=\text{ZF}\vee\text{SF} \neq \text{OF}$
\textbf{CC (Signed) x86}: \texttt{cmpq SRC, DST} then: \texttt{je / jz (DST == SRC),jne / jnz (DST != SRC),jl[e] (DST <[=]  SRC),jg[e] (DST >[=]  SRC)}
\textbf{Unsigned}: \texttt{je/jz (DST == SRC),jne/jnz (DST != SRC),jb[e] (DST <[=] SRC),ja[e] (DST >[=] SRC)}\\
\textbf{LLVM}: 
\texttt{\%c = icmp sgt i64 \%a \%b}. \texttt{\#c} is 1 iff \texttt{\%a<\%b}. More: \texttt{eq, ne, sgt, sge, slt, sle} 


\mysubsection{2.2 LLVM}
\textbf{Prefix}: local (\texttt{\%}) / global (\texttt{@}) identifier\\
\textbf{Phi-Nodes}: \lstinline[language=llvm]{%val = phi <ty> [%v1, %l1], [%v2, %l2]}
, select based most recent block visited\\ 
\textbf{GEP}:  \lstinline[language=llvm]{<res> = getelementptr <ty>, <ty>* <ptr>, <idxs>...},
purely arithmetic, if nested pointers \lstinline[language=llvm]{GEP(load(GEP(p*)))} \\


\textbf{Alloca \& Stack Storage:}  Allocates stack space, returns pointer, is freed on function return.

\textbf{Memory}: $8$(long), $4$(int) $\cdot \# $ vars. needed after func. call.
\begin{tabular}{@{}l l l@{}}
\toprule
\textbf{LLVM} & \textbf{x86 Equivalent} \\
\midrule
\texttt{sub i64 O1, O2} \textcolor{gray}{= o1 - o2} & \texttt{sub1 SRC, DST} \textcolor{gray}{= dst - src} \\
\texttt{\%L = alloca <ty>}  & \texttt{subq sizeof(T), \%rsp} \\
\texttt{\%L = load <ty>* P} & \texttt{movq (SRC), DST} \\
\texttt{store <ty> V, <ty>* P}  & \texttt{movq SRC, (DST)} \\
\midrule
\texttt{call <ty> F(...)} & \texttt{callq F} \\
\texttt{ret void / <ty>}  & \texttt{retq} \\
\texttt{br label \%L} & \texttt{jmp \%L} \\
\texttt{br i1 C, \%L1, \%L2} & \texttt{jne \%L1; jmp \%L2} \\
\midrule
\texttt{icmp <cond> ...} & \texttt{cmpq} + \texttt{jcc} \\
\texttt{getelementptr} & \texttt{leaq} (often) \\
\bottomrule
\end{tabular}






